\section{Bāzes prototipa analīze}
Šajā darbā izstrādātā mikrokontroliera kodols par paraugu izmanto
D.~Perija (\termEn{D.~Perry}) ,,VHDL: Programming by Example'' grāmatā
piedāvāto procesora imple\-men\-tā\-ciju \cite{Perry-VHDL}.

Sākotnēji (rev.~01) izstrādātais kodols izmantoja gandrīz identisku
arhitektūru (bet atšķirīgu instrukciju kopu), bet
jau agri izstrādes procesā tika identificētas vairāki arhitektūras trūkumi,
no kuriem daži padarīja prototipu nesintezējamu.
Šī nodaļa apskata šos trūkumus un piedāvātos risinājumus.
Problēmu risinājumi šajā nodaļā apskatīti visai virspusēji,
to implementācijas detaļas, pēc iestrādes kodolā,
sīkāk apskatītas \ref{sec:cpu}~un \ref{sec:uC}~nodaļā.

\begin{figure}[thb]
	\centering
	\includegraphics[scale=1.1]{perry-cpu}
	\caption[Bāzes prototipa uzbūves bloku diagramma.]
	        {Bāzes prototipa uzbūves bloku diagramma \cite[290.~lpp.]{Perry-VHDL}.}
	\label{fig:perry-cpu}
\end{figure}

D.~Perija procesors (turpmāk tekstā ,,bāzes prototips'')
ir vienkāršs 16 bitu procesors ar klasisku trīsstāvokļu 
centrālo datu kopni \cite{Flynn-arch}\cite{Heath}, un ar visām pamata
komponentēm (sk.~\ref{fig:perry-cpu}~att.), kas nodrošina bāzes prototipa
spēju izpildīt instrukcijas no tā instrukciju kopas.

\subsection{Iekšējā datu kopne} \label{sec:perry-bus}
\lipsum[18-18]

\subsection{Instrukciju kopa} \label{sec:perry-instr}
\lipsum[19-19]

\subsection{Komparators} \label{sec:perry-comp}
\lipsum[18-18]
